\subsection{Three reductions for four-object stars}

\begin{proposition}\label{C:fifteen-four-star-prunes}
Let $C$ be a four-object outward or inward star category with
at most fifteen arrows.  Write $M$ for its central endomorphism
monoid, $A_1,A_2,A_3$ for the leaf monoids, and $s_1,s_2,s_3$
for the nonempty leg cardinalities.  Each of the following
conditions excludes $BC\simeq\mathbb {CP}^2$:
\begin{enumerate}
\item Every $|A_i|\leq4$, and some $s_i=1$.
\item $|M|\leq9$ and every $|A_i|\leq2$.
\item At least two of $s_1,s_2,s_3$ equal one.
\end{enumerate}
\end{proposition}
\begin{proof}
Opposition interchanges inward and outward stars without changing
the classifying-space homotopy type.  For (1), take the inward
orientation.  The full subcategory $D$ on the leaves is a
disjoint union of three monoids of orders at most four, so
$H^2(BD;\mathbb Q)=0$.  A singleton incoming hom-set supplies
an arrow fixed by every element of $M$.
Lemma~\ref{C:fixed-incoming-cups}, which allows disconnected
$D$, makes all degree-two products zero.

For (2), take the outward orientation.  Every rational leaf
coinvariant functor is exact.  The actual orbit-collapse
functors replace the three leaves by trivial monoids and the
three legs by nonempty finite right $M$-sets $S_1,S_2,S_3$.
On each edge correspondence bar and each vertex they are
rational homology equivalences.  Hence they give a rational
cohomology-ring equivalence from $BC$ to the resulting star
$Z$.  Identify its three sink points.  Since no nonidentity
arrow exits a sink, the resulting nerve is exactly the act
cone for $S_1\sqcup S_2\sqcup S_3$.  Identifying three points
of a connected CW space gives the homotopy type
\[
 Z\vee S^1\vee S^1.
\]
If $BC\simeq\mathbb {CP}^2$, this act cone has bounded rational
homology, one-dimensional rational $H^2$, and a nonzero square.
Lemma~\ref{C:nine-bounded-rank-one} excludes this possibility;
the additional first homology is explicitly allowed there.

For (3), take the outward orientation and suppose $s_1=s_2=1$.
Put $X=BM$, and let $Y_i$ be the classifying space of the full
two-object subcategory on the center and leaf $i$.  The star
is the homotopy union of these three spaces along $X$; the
nerve inclusions are subcomplex inclusions.  A singleton leg
makes $Y_i$ the join $X*BA_i$.  Thus the maps $X\to Y_1$ and
$X\to Y_2$ are null-homotopic.  The ordinary homotopy-pushout
formula, first for these two maps and then for the third,
gives an actual homotopy equivalence
\begin{equation}\label{C:four-star-two-singleton-splitting}
 BC\simeq Y_1\vee Y_2\vee\Sigma X
       \vee\operatorname{hocofib}(X\longrightarrow Y_3).
\end{equation}
For the second pushout the common map from $X$ to
$Y_1\cup_X^hY_2$ is still null, since it factors through
the null map to $Y_1$.  This justifies the displayed splitting
without a choice of compatible inverse equivalences.

Suppose the left side were a projective plane.  The reduced
integral homology of the wedge on the right is the direct sum
of that of its four summands.  Its free-product fundamental
group is trivial, so each summand is simply connected.  The
unique degree-two integral generator and its primitive
nonzero square must belong to the same wedge summand; that
summand contains both nonzero homology groups, in degrees
two and four.  A suspension has zero positive cup products,
so $\Sigma X$ cannot be that summand.  It is therefore integrally
acyclic.  Since $X$ is connected, suspension homology shows
that $X$ is integrally acyclic as well.

Theorem~\ref{A:bounded} now makes the finite monoid space
$X=BM$ simply connected.  The homological Whitehead theorem
gives $X\simeq *$.  Consequently the two joins $Y_1,Y_2$ are
contractible and
$\operatorname{hocofib}(X\to Y_3)\simeq Y_3$.
Equation~\ref{C:four-star-two-singleton-splitting} reduces the
whole category to its center--third-leaf subcategory.  Its
number of arrows is at most
\[
 15-(|A_1|+1)-(|A_2|+1)\leq11.
\]
The already proved smaller-category exclusion is a contradiction.
This proves (3) and completes all three reductions.
\end{proof}
